\section{CREATE}
\subsection{Reward Channels and Agent State}
Let $R_t$ denote the generated reward program at iteration $t$. A policy $\pi_t$ is trained for a fixed interaction budget $B$ using only the generated reward $r_{R_t}$. A separate evaluator uses the unmodified environment reward $r_{\rm env}$:
\begin{equation}
\pi_t=\operatorname{PPO}(\mathcal{E},R_t;B),\qquad
J_t=\frac{1}{N}\sum_{i=1}^{N}\sum_{k=0}^{H_i-1}r_{\rm env}^{(i,k)}.
\label{eq:channels}
\end{equation}
This separation prevents the reward-design loop from evaluating a reward with its own generated signal.

The persistent agent state is
\begin{equation}
X_t=\{R_t,J_t,T_t,M_{t-1},A^*_{t-1}\},
\end{equation}
where $T_t$ contains episode length, termination, runtime errors, and named reward-component statistics; $M_{t-1}$ is the revision history; and $A^*_{t-1}$ is the best archived reward-policy pair. For a component $j$, the summary includes its nonzero activation rate and absolute magnitude share. These measurements are treated as diagnostic evidence, not as causal proof.

\begin{figure*}[t]
\centering
\includegraphics[width=0.98\textwidth]{create_framework.pdf}
\caption{CREATE places an autonomous reward-design loop above the policy learner. The generated reward is used only for PPO training; the environment-native reward is used only for evaluation. Training evidence and native-task performance drive the next reward revision, while lineage memory and a best archive preserve the evolution process.}
\label{fig:framework}
\end{figure*}

\subsection{Feedback-Guided Reward Revision}
The agent converts the state into a compact feedback record $F_t$ containing the threshold gap, termination pattern, component activation and scale, the previous revision, and the observed score change. Table~\ref{tab:evidence} summarizes the information contract. Raw policy trajectories and the native reward implementation are not exposed to the LLM; the agent receives aggregated measurements and the generated reward source. This keeps the design decision grounded in observed training outcomes while preserving the benchmark objective as an external evaluator.

\begin{table}[t]
\caption{Information contract of the reward-design agent.}
\label{tab:evidence}
\centering
\scriptsize
\setlength{\tabcolsep}{2.4pt}
\begin{tabular}{p{1.30cm}p{3.15cm}p{2.55cm}}
\toprule
Source & Agent-visible fields & Decision role \\
\midrule
Training & Component activation/scale, episode length, termination, runtime errors & Detect inactive, dominant, or invalid terms \\
Native eval. & Mean task return, range, threshold gap & Measure control-task progress \\
Lineage & Previous reward, declared revision, score response & Avoid repeated or destructive edits \\
Archive & Best reward and best native score & Protect the best discovered controller \\
\bottomrule
\end{tabular}
\end{table}

The update can be viewed as a supervisory control decision: the lower-level PPO learner maps observations to actions, whereas CREATE maps accumulated training evidence to a new learning objective. It therefore changes the policy-learning problem only between completed training runs, not the environment dynamics or the policy action at each time step.

It then produces the next reward:
\begin{equation}
R_{t+1}=\mathcal{U}(R_t,F_t,M_{t-1}).
\label{eq:update}
\end{equation}
Ordinary revisions modify one declared target: coefficients or gates are tuned when the structure remains usable, while one component is refactored when its mathematical form provides an ineffective signal. A global redesign is permitted only after repeated local failure. This bounded update preserves an interpretable lineage without claiming that a single edit identifies a counterfactual cause.

The archive stores the best development result,
\begin{equation}
A_t^*=\arg\max_{1\leq q\leq t}J_q,
\end{equation}
and early stopping is applied after the configured task threshold is reached or when further revisions cease to improve the active lineage. Algorithm~\ref{alg:create} summarizes the loop.

\begin{algorithm}[t]
\caption{CREATE Reward Evolution}
\label{alg:create}
\begin{algorithmic}[1]
\REQUIRE Environment $\mathcal{E}$, interface $I$, budget $B$, maximum rounds $K$, threshold $\theta$
\STATE $R_1\leftarrow\textsc{Generate}(I)$; $M_0\leftarrow\emptyset$; $A_0^*\leftarrow\emptyset$
\FOR{$t=1,\ldots,K$}
\STATE $\pi_t,T_t\leftarrow\textsc{Train}(\mathcal{E},R_t,B)$
\STATE $J_t,D_t\leftarrow\textsc{NativeEvaluate}(\mathcal{E},\pi_t)$
\STATE $A_t^*\leftarrow\textsc{ArchiveBest}(A_{t-1}^*,R_t,\pi_t,J_t)$
\IF{$J_t\geq\theta$ and stopping rule holds} \STATE \textbf{break} \ENDIF
\STATE $F_t\leftarrow\textsc{Summarize}(T_t,D_t,M_{t-1})$
\STATE $R_{t+1}\leftarrow\textsc{Revise}(R_t,F_t,M_{t-1})$
\STATE $M_t\leftarrow\textsc{Record}(M_{t-1},R_t,F_t,R_{t+1},J_t)$
\ENDFOR
\RETURN $A_t^*$
\end{algorithmic}
\end{algorithm}
